% GENERATED FILE. Edit canonical lessons, then run npm run generate:books.
% canonical-source-hash: fnv1a64:997b3de5161b06eb
% canonical-lessons: BN-C326-ashchoryo, BN-C326-oloukik, BN-C326-kolpona-kora, BN-C326-guptodhon, BN-C326-miliye-jaoya

\chapter{Wonders}
\label{ch:bn-wonders}
\hlchaptermodality{326}

\begin{chapteropening}
\textbf{By the end of this chapter:} \emph{I can describe wonders, miracles and hidden treasure.}

Complete the last lesson of chapter 326: I can describe wonders, miracles and hidden treasure.
\end{chapteropening}

\section[\texorpdfstring{\bn{আশ্চর্য} (āshchôryô)}{আশ্চর্য (āshchôryô)}]{\bn{আশ্চর্য} (āshchôryô) — amazing; a wonder}

\label{lesson:BN-C326-ashchoryo}



\begin{quote}
Before the new one: say the Bengali for to mend one's ways, then the Bengali for to be forced.
\end{quote}

\subsection*{You'll want to know: \bn{আশ্চর্য}}
\textbf{\bn{আশ্চর্য}} — \emph{āshchôryô} — \textquotedblleft{}amazing; a wonder\textquotedblright{}.

\begin{quote}
\bn{সন্ন্যাসীর} \bn{গল্পটা} \bn{সত্যি} \bn{আশ্চর্য} \bn{ছিল।}
\end{quote}

\begin{quote}
\emph{sônnyāsīr gôlpoṭā sôtti āshchôryô chhilô.}
\end{quote}

\begin{quote}
The holy man's story was truly amazing.
\end{quote}

\begin{cousinweb}
\textbf{\bn{আশ্চর্য}} (\emph{āshchôryô}) is Sanskrit \emph{āścarya}, 'appearing rarely, wonderful; a strange sight' (Monier-Williams).
\end{cousinweb}

\begin{grammarlens}[title={What you've built}]
One more word for telling what happened.
\end{grammarlens}

\subsection*{Your turn}
\begin{itemize}
\raggedright
  \item \emph{Say it:} \emph{āshchôryô}
  \item \emph{Say it:} \emph{āshchôryô}, once more
  \item \emph{Say it:} \emph{sônnyāsīr gôlpoṭā sôtti āshchôryô chhilô.} — The holy man's story was truly amazing.
  \item \emph{Recall:} say the Bengali for to mend one's ways, then the Bengali for to be forced, then say \emph{āshchôryô} again
\end{itemize}

\begin{tcolorbox}[breakable,colback=teal!4,colframe=teal!35!black,title={Before you move on}]
What does \bn{আশ্চর্য} mean? (\textquotedblleft{}Amazing; a wonder\textquotedblright{}.) Say it once more.
\end{tcolorbox}
\section[\texorpdfstring{\bn{অলৌকিক} (ôloukik)}{অলৌকিক (ôloukik)}]{\bn{অলৌকিক} (ôloukik) — supernatural; miraculous}

\label{lesson:BN-C326-oloukik}



\begin{quote}
Before the new one: say the Bengali for to be forced, then the Bengali for amazing.
\end{quote}

\subsection*{You'll want to know: \bn{অলৌকিক}}
\textbf{\bn{অলৌকিক}} — \emph{ôloukik} — \textquotedblleft{}supernatural; miraculous\textquotedblright{}.

\begin{quote}
\bn{লোকে} \bn{বলে} \bn{যে} \bn{মন্দিরের} \bn{পুকুরের} \bn{জলে} \bn{অলৌকিক} \bn{শক্তি} \bn{আছে।}
\end{quote}

\begin{quote}
\emph{loke bôle je mondirer pukurer jôle ôloukik shôkti āchhe.}
\end{quote}

\begin{quote}
People say the water of the temple pond has supernatural power.
\end{quote}

\begin{cousinweb}
\textbf{\bn{অলৌকিক}} (\emph{ôloukik}) is Sanskrit \emph{alaukika}, 'uncommon; not of this world, supernatural' (Monier-Williams): \emph{a-}, 'not', and \emph{laukika}, 'worldly'.
\end{cousinweb}

\begin{grammarlens}[title={What you've built}]
One more word for telling what happened.
\end{grammarlens}

\subsection*{Your turn}
\begin{itemize}
\raggedright
  \item \emph{Say it:} \emph{ôloukik}
  \item \emph{Say it:} \emph{ôloukik}, once more
  \item \emph{Say it:} \emph{loke bôle je mondirer pukurer jôle ôloukik shôkti āchhe.} — People say the water of the temple pond has supernatural power.
  \item \emph{Recall:} say the Bengali for to be forced, then the Bengali for amazing, then say \emph{ôloukik} again
\end{itemize}

\begin{tcolorbox}[breakable,colback=teal!4,colframe=teal!35!black,title={Before you move on}]
What does \bn{অলৌকিক} mean? (\textquotedblleft{}Supernatural; miraculous\textquotedblright{}.) Say it once more.
\end{tcolorbox}
\section[\texorpdfstring{\bn{কল্পনা} \bn{করা} (kôlpônā kôrā)}{কল্পনা করা (kôlpônā kôrā)}]{\bn{কল্পনা} \bn{করা} (kôlpônā kôrā) — to imagine}

\label{lesson:BN-C326-kolpona-kora}



\begin{quote}
Before the new one: say the Bengali for amazing, then the Bengali for supernatural.
\end{quote}

\subsection*{You'll want to know: \bn{কল্পনা} \bn{করা}}
\textbf{\bn{কল্পনা} \bn{করা}} — \emph{kôlpônā kôrā} — \textquotedblleft{}to imagine\textquotedblright{}.

\begin{quote}
\bn{রাতে} \bn{আমি} \bn{জাদুর} \bn{দেশ} \bn{কল্পনা} \bn{করতাম।}
\end{quote}

\begin{quote}
\emph{rāte āmi jādur desh kôlpônā kôrtām.}
\end{quote}

\begin{quote}
At night I used to imagine a land of magic.
\end{quote}

Past (he/she …): \bn{তিনি} \bn{একটা} \bn{নতুন} \bn{শহর} \bn{কল্পনা} \bn{করলেন}, \emph{tini ekṭā nôtun shôhôr kôlpônā kôrlen}.

\begin{cousinweb}
\textbf{\bn{কল্পনা} \bn{করা}} (\emph{kôlpônā kôrā}) puts Sanskrit \emph{kalpanā}, 'making, forming' (Monier-Williams), before \bn{করা}: to form a picture in the mind.
\end{cousinweb}

\begin{grammarlens}[title={What you've built}]
One more word for telling what happened.
\end{grammarlens}

\subsection*{Your turn}
\begin{itemize}
\raggedright
  \item \emph{Say it:} \emph{kôlpônā kôrā}
  \item \emph{Say it:} \emph{kôlpônā kôrā}, once more
  \item \emph{Say it:} \emph{rāte āmi jādur desh kôlpônā kôrtām.} — At night I used to imagine a land of magic.
  \item \emph{Recall:} say the Bengali for amazing, then the Bengali for supernatural, then say \emph{kôlpônā kôrā} again
\end{itemize}

\begin{tcolorbox}[breakable,colback=teal!4,colframe=teal!35!black,title={Before you move on}]
What does \bn{কল্পনা} \bn{করা} mean? (\textquotedblleft{}To imagine\textquotedblright{}.) Say it once more.
\end{tcolorbox}
\section[\texorpdfstring{\bn{গুপ্তধন} (guptôdhôn)}{গুপ্তধন (guptôdhôn)}]{\bn{গুপ্তধন} (guptôdhôn) — a hidden treasure}

\label{lesson:BN-C326-guptodhon}



\begin{quote}
Before the new one: say the Bengali for supernatural, then the Bengali for to imagine.
\end{quote}

\subsection*{You'll want to know: \bn{গুপ্তধন}}
\textbf{\bn{গুপ্তধন}} — \emph{guptôdhôn} — \textquotedblleft{}a hidden treasure\textquotedblright{}.

\begin{quote}
\bn{গল্পে} \bn{নায়ক} \bn{একটা} \bn{পুরনো} \bn{কুয়োর} \bn{নিচে} \bn{গুপ্তধন} \bn{পেল।}
\end{quote}

\begin{quote}
\emph{gôlpe nāyôk ekṭā purôno kuyor niche guptôdhôn pelô.}
\end{quote}

\begin{quote}
In the story the hero found hidden treasure at the bottom of an old well.
\end{quote}

\begin{cousinweb}
\textbf{\bn{গুপ্তধন}} (\emph{guptôdhôn}) joins Sanskrit \emph{gupta}, 'guarded; hidden' (Monier-Williams), to \emph{dhana}, 'wealth': hidden wealth.
\end{cousinweb}

\begin{grammarlens}[title={What you've built}]
One more word for telling what happened.
\end{grammarlens}

\subsection*{Your turn}
\begin{itemize}
\raggedright
  \item \emph{Say it:} \emph{guptôdhôn}
  \item \emph{Say it:} \emph{guptôdhôn}, once more
  \item \emph{Say it:} \emph{gôlpe nāyôk ekṭā purôno kuyor niche guptôdhôn pelô.} — In the story the hero found hidden treasure at the bottom of an old well.
  \item \emph{Recall:} say the Bengali for supernatural, then the Bengali for to imagine, then say \emph{guptôdhôn} again
\end{itemize}

\begin{tcolorbox}[breakable,colback=teal!4,colframe=teal!35!black,title={Before you move on}]
What does \bn{গুপ্তধন} mean? (\textquotedblleft{}A hidden treasure\textquotedblright{}.) Say it once more.
\end{tcolorbox}
\section[\texorpdfstring{\bn{মিলিয়ে} \bn{যাওয়া} (miliye jāoyā)}{মিলিয়ে যাওয়া (miliye jāoyā)}]{\bn{মিলিয়ে} \bn{যাওয়া} (miliye jāoyā) — to vanish; to fade away}

\label{lesson:BN-C326-miliye-jaoya}



\begin{quote}
Before the new one: say the Bengali for to imagine, then the Bengali for a hidden treasure.
\end{quote}

\subsection*{You'll want to know: \bn{মিলিয়ে} \bn{যাওয়া}}
\textbf{\bn{মিলিয়ে} \bn{যাওয়া}} — \emph{miliye jāoyā} — \textquotedblleft{}to vanish; to fade away\textquotedblright{}.

\begin{quote}
\bn{এক} \bn{মুহূর্তে} \bn{পরীটা} \bn{বাতাসে} \bn{মিলিয়ে} \bn{গেল।}
\end{quote}

\begin{quote}
\emph{ek muhūrte pôrīṭā bātāse miliye gelô.}
\end{quote}

\begin{quote}
In a moment the fairy vanished into thin air.
\end{quote}

Past (he/she …): \bn{সেই} \bn{আলো} \bn{মিলিয়ে} \bn{গেল}, \emph{sei ālo miliye gelô}.

\begin{cousinweb}
\textbf{\bn{মিলিয়ে} \bn{যাওয়া}} (\emph{miliye jāoyā}) is the -e form of \bn{মিলানো}, to blend in — CDIAL 10133, Sanskrit \emph{milati}, 'meets', which lists Bengali \emph{milā} — before \textbf{\bn{যাওয়া}}: to blend into the air and be gone.
\end{cousinweb}

\begin{grammarlens}[title={What you've built}]
One more word for telling what happened.
\end{grammarlens}

\subsection*{Your turn}
\begin{itemize}
\raggedright
  \item \emph{Say it:} \emph{miliye jāoyā}
  \item \emph{Say it:} \emph{miliye jāoyā}, once more
  \item \emph{Say it:} \emph{ek muhūrte pôrīṭā bātāse miliye gelô.} — In a moment the fairy vanished into thin air.
  \item \emph{Recall:} say the Bengali for to imagine, then the Bengali for a hidden treasure, then say \emph{miliye jāoyā} again
\end{itemize}

\begin{tcolorbox}[breakable,colback=teal!4,colframe=teal!35!black,title={Before you move on}]
What does \bn{মিলিয়ে} \bn{যাওয়া} mean? (\textquotedblleft{}To vanish; to fade away\textquotedblright{}.) Say it once more.
\end{tcolorbox}
